\chapter{Conclusion}
\label{ch:conclusion}


\section{Summary of Findings}

This thesis asked one question: in real-time sepsis prediction on MIMIC-IV, how
much of the reported performance comes from the model versus the researcher's
choices about the label, the data split, and the evaluation metric?

The answer is not the one the study anticipated. Model class dominates
discrimination. External transport comes second, moving AUROC several times
more than the label definition does, though on an interval too wide to pin
down. The label and the split design move it least. What the label determines is
not how well a cohort can be predicted but \emph{which cohort exists at all},
and, through its onset rule, whether prediction ahead of clinical action can be
measured in the first place.

\paragraph{H1 (Label dominance): Not confirmed, and reversed.}
Changing the sepsis definition while keeping the model fixed changes AUROC by
\pcSpreadLabel. Changing the model while keeping the label fixed changes it by
\pcSpreadModelClass{}, more than tenfold, and the interval on the
difference, [\pcHOneCiLo, \pcHOneCiHi], excludes zero. Cohen et al.'s
\cite{cohen2024subtle} MIMIC-III finding does not reproduce on MIMIC-IV as a
\emph{discrimination} effect. It does reproduce as a cohort effect: the two
narrowest readings agree on which stays are septic at only
$\kappa = \pcKappaVsBA$ despite near-identical prevalence, so the label still
determines who is counted, it simply does not determine how learnable they
are.

\paragraph{H2 (Coefficient instability): Criterion met, descriptively.}
\pcHTwoEst{} of the \pcHTwoNTerms{} shared covariates are unstable across label
variants: \pcHTwoNSign{} reverse sign (bilirubin and vasopressor exposure) and
\pcHTwoNMag{} change magnitude at least twofold, three of the latter being
missingness indicators rather than physiological values. Vasopressor exposure
meets both criteria, which is why the two counts sum to seven across six
covariates. H2 is a stability criterion rather than a contrast, so no test
statistic is defined for it, and the diagnostic of
Section~\ref{sec:h2_uncertainty} shows the count to be a loose upper bound:
\pcHTwoNNoise{} of the \pcHTwoNFlagged{} flagged covariates, both sign
reversals among them, move by less than their own estimation noise. What
survives that comparison is confined to the order indicators, so the
defensible reading is narrower than the count: an effect estimate for whether
a test was ordered cannot be interpreted without knowing the exact label
definition.

\paragraph{H3 (Treatment leakage): Definitional, plus its empirical counterpart.}
The pre-treatment window contains \pcPreTreatOnsetsVarB{} sepsis onset events
across \pcPreTreatPersonHoursVarB{} person-hours, but that zero is entailed by
the onset rule the pre-registered variants use, not measured
(Proposition~\ref{prop:h3}), so the pre-registered contrast was never testable.
The empirical question was answered instead by relabelling the identical stays
under the standard $\min(t_{\text{susp}}, t_{\text{SOFA}})$ onset of Seymour et
al.\ and the PhysioNet/CinC 2019 Challenge (Variant F, Protocol Amendment 4):
\pcFNMovedEarlier{} onsets move earlier and \pcPreTreatOnsetsVarF{} fall before
treatment, identified at \pcPreTreatAurocPrimaryF{} AUROC. A published MIMIC-IV
real-time sepsis model measures anticipation of physiological deterioration only
to the extent its onset rule allows an onset to precede clinical action; under a
conjunction rule, it cannot.

\paragraph{H4 (Split optimism): Not confirmed.}
The random split does not flatter the model. For the primary model it scores
\pcAurocRandomPrimaryB{} against \pcAurocPrimaryB{} on the temporal split, a
difference of \pcHFourEst{} (95\% CI \pcHFourCiLo--\pcHFourCiHi) in the
direction opposite to the hypothesis, and every model fitted under both splits
shows the same sign. The two test sets differ in composition as well as in
design (the temporal arm is a single admission era, the random arm a sample
across all of them), so this is not a clean optimism estimate; what it does
establish is that a random split cannot be assumed to be the more permissive
choice.

\paragraph{H5 (Metric divergence): Not confirmed, and superseded by a starker
finding.}
Neither arm supports the hypothesis, but the AUROC arm fails narrowly and in an
informative way. External discrimination is \pcExtAurocPrimaryB{} on eICU-CRD
against \pcAurocPrimaryB{} internally, a degradation of \pcHFiveEst{}: the
predicted direction, and roughly three-quarters of the 0.085 anticipated from
Moor et al.\ \cite{moor2023review}. The interval (\pcHFiveCiLo--\pcHFiveCiHi)
nonetheless covers zero, because the Sepsis-3 anchor is recoverable in only
\pcExtMicroCoverage\% of eICU-CRD stays and the estimate rests on
\pcExtNEventsB{} events. This is the one hypothesis a better-powered external
cohort might have confirmed, and the antibiotic-only arm (\pcExtAbxNEventsB{}
events, a different target) shows the same degradation with the same ordering
across labels. The study establishes that discrimination declines on transport
and cannot establish by how much.

The utility arm needs no such qualification, because it is settled internally.
The proportional decline the hypothesis asked for cannot be formed at all,
there being no internal utility to decline from: swept across sixty thresholds,
the best any configuration achieves under the primary label is
\pcUtilityMaxGbtB{}, at \pcUtilityOptBurdenGbtB{} false alerts per true one
against a benchmark of 1.4. The utility collapse Wang et al.\
\cite{wang2025srpm} documented across ninety-one models therefore has its
counterpart \emph{internally} on MIMIC-IV under an honest time-based split,
there was scarcely any margin left to collapse. Their broader point survives in
strengthened form: on the cohort where the label is fully observed, AUROC and
clinical utility are close to uncorrelated, discrimination sitting in the high
0.7s while usable clinical value never materialises at any threshold.

\paragraph{H6 (Model-class null): Not rejected; non-inferiority established.}
The interpretable hazard model is not outperformed by gradient-boosted trees:
\pcAurocPrimaryB{} against \pcAurocGbtB, a difference of \pcHSixEst{} with an
interval of [\pcHSixCiLo, \pcHSixCiHi]. The upper limit sits inside the
pre-registered 0.02 margin, so the trees are not better by more than it; the
lower limit falls outside, so two-sided equivalence is \emph{not} established
and a difference of up to roughly \pcHSixCiLo{} in the interpretable model's
favour remains consistent with these data. This is the strongest positive
result in the study, and it is strongest in exactly the direction that matters:
on tabular ICU time-series of this kind the accuracy case for preferring an
opaque model is not supported. What little AUROC separates them buys nothing
clinically: at their respective best thresholds the two models reach
\pcUtilityMaxGbtB{} and \pcUtilityMaxPrimaryB{} on the Utility Score, a gap of
little more than a hundredth of a point on a scale bounded by 1, and both sit
within a few per cent of the no-alert strategy. Discrimination and clinical
value part company here not by re-ordering the models but by disagreeing about
whether the difference matters.


\section{The Central Lesson}

The utility collapse documented by Wang et al.\ \cite{wang2025srpm} is
not a model problem. It is a label problem. The Sepsis-3 label on MIMIC-IV:

\begin{enumerate}
    \item \textbf{Depends on clinical action for membership:} no patient is
    \emph{labelled} septic without a qualifying antibiotic-culture pair, under
    any timing rule. The further claim that no onset can \emph{precede} that
    pair is a property of the conjunction-based onset rule alone: timed as the
    standard specifies, Sepsis-3 records \pcPreTreatOnsetsVarF{} pre-treatment
    onsets on the same stays, and a treatment-independent deterioration label
    records \pcPreTreatOnsetsVarD{}. The emptiness belonged to the
    implementation, not to the patients.
    \item \textbf{Is sensitive to how it is defined:} the number of sepsis
    cases nearly doubles across plausible Sepsis-3 interpretations, and the two
    narrowest readings agree on which stays are septic at only
    $\kappa = \pcKappaVsBA$. Discrimination is largely insensitive to this
    (\pcSpreadLabel{} of AUROC); every count, rate and denominator built on the
    label is not.
    \item \textbf{Produces unusable alert behaviour:} no model, under any label
    variant, has a threshold at which it both beats the no-alert strategy by a
    clinically meaningful margin and raises an alarm volume a ward could
    absorb. The best configuration under the primary label reaches
    \pcUtilityMaxGbtB{} at \pcUtilityOptBurdenGbtB{} false alerts per true one.
\end{enumerate}

These three problems are independent of model complexity. A better model
cannot fix a label that is constituted by the very clinical actions it is
supposed to predict ahead of. The solution requires changing the label.

Splitting the definition sharpens point (2) and qualifies it
(Section~\ref{sec:label_anchor_results}). It is qualified because the onset
\emph{time} under every pre-registered Sepsis-3 variant is fixed entirely by the
treatment anchor, the organ-dysfunction criterion decides which stays are
labelled and moves no onset by an hour, so part of the variation between
variants is guaranteed by the definition rather than found in the data. It is
sharpened because a label built from physiology alone, owing nothing to
treatment timing, agrees with Sepsis-3 barely better than chance
($\kappa = \pcKappaVsBD$). The label effect is therefore real, not merely
definitional.

The same comparison relocates the difficulty, and not where the conjunction
framing expected. The treatment anchor on its own is as predictable as Sepsis-3
built on top of it (the anchor-attributable fraction is
\pcPhiAnchorPrimary{} with an interval
[\pcPhiAnchorCiLoPrimary, \pcPhiAnchorCiHiPrimary] that includes 1), so
requiring both limbs is not what makes the label hard. Physiological
deterioration alone is \emph{easier} (\pcAltAurocFullPrimaryD{} against
\pcAurocPrimaryB), and Sepsis-3 with a better clock is easier again
(\pcAltAurocFullPrimaryF). What resists prediction is not the conjunction but
the moment at which the conjunction is stamped: the clinician's action rather
than the patient's deterioration.


\section{Contributions}

\begin{enumerate}
    \item \textbf{C1 (Label variance, separated from discrimination):} Three
    Sepsis-3 readings applied to a discrete-time competing-risks model on
    MIMIC-IV (\pcNStays{} stays, against the roughly two thousand admissions of
    Cohen et al.'s MIMIC-III comparison), with the label and model-class
    spreads estimated on the same test rows and each carrying an interval. That
    combination is, to the best of our knowledge, not previously reported, and
    it is what separates two effects the literature treats as one: label choice
    moves cohort membership substantially ($\kappa = \pcKappaVsBA$ between the
    narrowest readings) while moving AUROC by only \pcSpreadLabel.

    \item \textbf{C2 (Treatment leakage located in the onset rule):} A proof that
    the conjunction-based Sepsis-3 onset rule makes the pre-treatment window
    empty by construction, paired with the measurement that the standard
    $\min(t_{\text{susp}}, t_{\text{SOFA}})$ rule recovers
    \pcPreTreatOnsetsVarF{} pre-treatment onsets on the identical stays at
    \pcPreTreatAurocPrimaryF{} AUROC. This confirms Kamran et al.'s
    \cite{kamran2024evaluation} concern for label membership and corrects its
    over-extension to label timing.

    \item \textbf{C3 (Utility ceiling):} A threshold-swept Utility Score for
    every model and label variant, establishing that the ceiling, not merely
    the score at a conventional cut-off, sits within a few per cent of the
    no-alert strategy under the primary label, and that reaching even that
    ceiling costs
    \pcUtilityOptBurdenGbtB{} false alerts per true one. This extends Wang et
    al.'s \cite{wang2025srpm} meta-analytic finding to a controlled
    single-study comparison and removes the operating-point objection that
    would otherwise apply to it.

    \item \textbf{C4 (Coefficient instability):} \pcHTwoEst{} of the
    \pcHTwoNTerms{} shared covariates are unstable across labels
    (\pcHTwoNSign{} by sign and \pcHTwoNMag{} by magnitude, vasopressor
    exposure meeting both), supporting Lauritsen et al.'s
    \cite{lauritsen2021framing} framing-instability claim on MIMIC-IV for
    specific effect estimates, measured against an identified fit. Referred to
    their own cluster-robust standard errors, \pcHTwoNRobust{} of the
    \pcHTwoNFlagged{} survive, and all \pcHTwoNRobust{} are order indicators
    (Section~\ref{sec:h2_uncertainty}).

    \item \textbf{C5 (Equity-by-label):} Subgroup discrimination and subgroup
    label sensitivity measured on the same strata in the same
    onset-prediction cohort, both with cluster-bootstrap intervals under
    false-discovery control: a pairing this study's search did not find
    reported elsewhere, though each component separately has precedent
    (Section~\ref{sec:search_strategy}). The label-sensitivity contrasts
    survive correction and the discrimination contrasts do not, so the
    defensible finding is about the completeness of the demographic record
    rather than about race; reporting either analysis without the other would
    have inverted which conclusion the data support.

    \item \textbf{C6 (Reproducible pipeline):} A pre-registered,
    TRIPOD+AI-conformant pipeline with fixed seeds and a single orchestration
    script, which reproduces every reported estimate exactly from the raw
    databases (with the one documented exception of the gradient-boosted
    comparator, whose fitting is not bit-reproducible across library builds and
    whose rows therefore serve as the designated consistency check on any
    re-run), and which emits the numbers in this document as generated
    constants so that text and code cannot diverge
    (Section~\ref{sec:code_availability}, Appendix~\ref{app:code}).

    \item \textbf{C7 (Definitional vs.\ empirical label variance separated):}
    A treatment-independent deterioration label, tested alongside the three
    Sepsis-3 variants, shows that the label effect is not an artefact of how one
    definition is read: it agrees with Sepsis-3 at $\kappa = \pcKappaVsBD$ and
    matches its onset hour in \pcPctExactMatchD\% of shared stays. The same
    comparison shows that the treatment anchor alone is as predictable as
    Sepsis-3 built on top of it ($\phi_{\text{anchor}} = \pcPhiAnchorPrimary$,
    interval including 1), and that the pre-treatment window empty of Sepsis-3
    onsets contains
    \pcPreTreatOnsetsVarD{} treatment-independent deterioration onsets
    identified at \pcPreTreatAurocPrimaryD{} AUROC, and
    \pcPreTreatOnsetsVarF{} Sepsis-3 onsets once the standard onset timing is
    restored. This relocates the difficulty of sepsis prediction from the
    physiology to the definition, and within the definition to the onset rule
    (Section~\ref{sec:label_anchor_results}).
\end{enumerate}


\section{Future Work}

\begin{enumerate}
    \item \textbf{Better sepsis labels.} Section~\ref{sec:label_anchor_results}
    takes the first step (a deterioration label built from physiology alone,
    with no treatment timestamp), and shows it identifies substantially
    different patients at substantially different hours. But that label is an
    operational construct built for this comparison, not a validated definition,
    and its components overlap the model's own inputs. The work remaining is a
    sepsis definition grounded in objective measurements such as lactate
    trajectories and organ-function biomarkers, adjudicated independently of the
    model's feature set, which would eliminate the treatment-constitutionality
    problem and allow genuine pre-onset prediction.

    \item \textbf{Non-Sepsis-3 consensus definitions.} SEP-1 and the CDC Adult
    Sepsis Event were not implemented. Neither would have tested
    treatment-independence, both are defined on antibiotic days, but each
    yields a different cohort, and Dutta et al.\ \cite{dutta2026performance}
    report that performance moves substantially between them. Adding them would
    extend the label-variance decomposition across definitions rather than
    within one.

    \item \textbf{Prospective validation.} Test the new label in real clinical
    settings where sepsis onset can be determined by a protocol-blinded assessor,
    independent of treatment decisions.

    \item \textbf{Utility-optimised training.} The thresholds reported here are
    optimised post hoc, but the models themselves are not: each is fitted to a
    likelihood or a squared-error surrogate and only afterwards read at the
    threshold that happens to maximise utility. Training directly against the
    Utility Score (as a differentiable surrogate loss, or through a
    cost-sensitive reweighting of the person-hour rows) would change what the
    model learns rather than only how it is read, and is the natural next step
    given how shallow the swept optima proved to be. The gap between such a
    model and the post-hoc ceiling in Table~\ref{tab:utility_optimum} would also
    quantify how much of the failure is representational rather than
    operational.

    \item \textbf{Multi-hospital evaluation.} The alert burden, \pcUtilityOptBurdenGbtB{} false alerts per true alert even at the best
    model's best operating point, is a deployment problem, not a statistical
    one. Determining whether any model can be deployed at a practically
    acceptable burden requires prospective multi-centre testing with
    clinical-staff feedback.
\end{enumerate}


\section{Closing Remark}

The premise this study was built on did not survive the study. It was designed
around the expectation that the analyst's discretionary choices would rival the
choice of model as a source of performance variation, and that expectation
fails: the label moves discrimination by \pcSpreadLabel{} against the model
class's \pcSpreadModelClass{}, and the split method moves it not at all. Of the
four hypotheses that could be tested, none is confirmed, and two fail in the
direction opposite to the one predicted. That outcome is reported here in the
same detail it would have received had it gone the other way, which is the
reason the pre-registration and the inference layer were worth their cost: they
fix what the answer means before the answer is known, and on point estimates
alone this thesis would have claimed the opposite of what it can defend.

The premise inverted; the motivating problem did not. What the decomposition
establishes is that the label's consequence does not run through
discrimination, which is where the field looks for it. It runs through cohort
membership, the two narrowest readings agreeing at $\kappa = \pcKappaVsBA$ on
which patients are septic, and through onset timing, which fixes what a
real-time model can be asked to anticipate at all. Those are the channels the reported
metric is blind to, which is why a literature that has optimised AUROC for a
decade has not converged on models that help.

The central finding, that the Sepsis-3 label on MIMIC-IV is constituted by
clinical action, is not a failure of the model. It is a finding about the label
itself, and it is the most important contribution this study makes.

The last part of that finding required looking outside the pre-registered
variants. The pre-treatment window is not empty of predictable deterioration,
it contains \pcPreTreatOnsetsVarD{} onsets under a treatment-independent label
at \pcPreTreatAurocPrimaryD{} AUROC, and \pcPreTreatOnsetsVarF{} onsets under
Sepsis-3 itself once the onset is timed as
$\min(t_{\text{susp}}, t_{\text{SOFA}})$ rather than as the suspicion time
alone. It was empty here because of an onset rule, and that rule is not
peculiar to this thesis. The signal clinicians would want a model to find is
present in the data; a widely used way of writing the label down is what makes
it unrepresentable.
